\documentclass[a4paper]{article}
\usepackage[fontset=fandol]{ctex}
\usepackage{amsmath, amssymb}
\usepackage{graphicx}
\usepackage[margin=2.5cm]{geometry}
\usepackage[most]{tcolorbox}
\usepackage{etoolbox}
\usepackage{listings}
\usepackage{hyperref}
\usepackage{booktabs}
\usepackage{subcaption}
\usepackage{float}
\usepackage{tikz}
\IfFileExists{pgfplots.sty}{\usepackage{pgfplots}\pgfplotsset{compat=1.18}}{}
\newtcolorbox{knowledgebox}[1]{enhanced, colback=blue!5!white, colframe=blue!75!black, colbacktitle=blue!75!black, coltitle=white, fonttitle=\bfseries, title=#1, attach boxed title to top left={yshift=-2mm, xshift=2mm}, boxrule=1pt, sharp corners}
\newtcolorbox{importantbox}[1]{enhanced, colback=yellow!10!white, colframe=yellow!80!black, colbacktitle=yellow!80!black, coltitle=black, fonttitle=\bfseries, title=#1, sharp corners}
\newtcolorbox{warningbox}[1]{enhanced, colback=red!5!white, colframe=red!75!black, colbacktitle=red!75!black, coltitle=white, fonttitle=\bfseries, title=#1, sharp corners}
\lstset{language=Python,basicstyle=\ttfamily\small,keywordstyle=\color{blue},stringstyle=\color{red!60!black},commentstyle=\color{green!60!black},breaklines=true,frame=single,numbers=left,numberstyle=\tiny\color{gray},captionpos=b,extendedchars=false}
\newcommand{\notetitle}{Guest Lecture II}
\newcommand{\noteauthors}{Codex 基于课程页、公开视频字幕、coverage ledger 与补充官方材料整理}
\newcommand{\notedate}{\today}
\newcommand{\videochannel}{WAVLab}
\newcommand{\videopublishdate}{2023-11-29}
\newcommand{\videoduration}{unknown}
\newcommand{\videourl}{}
\newcommand{\videocoverpath}{}
\begin{document}
\begin{titlepage}
\centering
{\Large 课程笔记\par}
\vspace{1.2cm}
{\huge\bfseries \notetitle\par}
\vspace{0.8cm}
{\large \noteauthors\par}
\vspace{0.3cm}
{\large \notedate\par}
\vspace{1.2cm}
\ifdefempty{\videocoverpath}{{\small 公开视频封面不可用。\par}}{\includegraphics[width=0.82\textwidth,height=0.45\textheight,keepaspectratio]{\videocoverpath}\par}
\vfill
\begin{tcolorbox}[width=0.9\textwidth, colback=black!2!white, colframe=black!60, sharp corners]
\textbf{视频作者/频道}：\videochannel\par
\textbf{发布日期}：\videopublishdate\par
\textbf{视频时长}：\videoduration\par
\textbf{视频链接}：\href{\videourl}{\nolinkurl{\videourl}}
\end{tcolorbox}
\end{titlepage}
\tableofcontents
\newpage
\begin{warningbox}{Source Discrepancy}
2023 course page exposes only a Guest Lecture placeholder. This chapter is reconstructed from the official 2023 schedule row, the official same-course 2022 guest lecture metadata naming Thomas Shaf at 3M | MModal, and official 3M / Solventum product documentation.
\end{warningbox}

\section{23.0 重建前言：从课程排程到工业语音系统}
根据课程页主题、公开视频字幕以及本讲的 coverage ledger，23.0 重建前言：从课程排程到工业语音系统 这一部分主要围绕 2023 course row: Guest Lecture | no public title/video/slide on the 2023 page. 展开。
写作上保持 coverage-first 原则：先把老师在这一阶段真正讲过的问题设定、模型部件、训练或推理逻辑讲清，再压缩成适合复习的结构。

\begin{figure}[H]
\centering
\includegraphics[width=0.92\textwidth]{figures/coverage_map.png}
\caption{本讲 coverage map：按 coverage\_units.jsonl 归纳出的核心主题与章节映射。}
\end{figure}

\begin{figure}[H]
\centering
\includegraphics[width=0.92\textwidth]{figures/timeline.png}
\caption{本讲 timeline：根据字幕时间范围归纳的主题推进顺序。}
\end{figure}

\subsection{2023 course row: Guest Lecture}
从证据位置看，这一部分对应的课程证据区间是 2023-11-29。课程在这里强调的核心内容可以概括为：2023 course row: Guest Lecture | no public title/video/slide on the 2023 page.。
从课程组织上看，这一部分承担的是把局部概念嵌回整条识别流水线中的作用，避免读者把公式或模块孤立地记忆。

记录系统中的附加说明指出：This chapter is an explicitly labeled reconstruction, not a recovered 2023 verbatim lecture.

在本章结构中，这一单元被并入“23.0 重建前言：从课程排程到工业语音系统”，目的是把同一阶段的定义、机制和权衡集中讲清。

\subsection{本章小结}
23.0 重建前言：从课程排程到工业语音系统 这一部分把 Guest Lecture II 中相邻的教学单元串成了一条连续叙事：既交代了这一阶段要解决的问题，也交代了方法、推理或工程权衡为什么要这样组织。

\section{23.1 Guest Lecture Reconstruction: Industrial Speech Systems and Clinical Documentation}
根据课程页主题、公开视频字幕以及本讲的 coverage ledger，23.1 Guest Lecture Reconstruction: Industrial Speech Systems and Clinical Documentation 这一部分主要围绕 同门课 2022 官方 guest lecture row 明确写出 Thomas Shaf at 3M | MModal，说明这一块是工业 speech system 主题。 展开。
写作上保持 coverage-first 原则：先把老师在这一阶段真正讲过的问题设定、模型部件、训练或推理逻辑讲清，再压缩成适合复习的结构。

\subsection{同门课 2022 官方 guest lecture row 明确写出 Thomas Shaf at 3M}
从证据位置看，这一部分对应的课程证据区间是 2022-11-30 row。课程在这里强调的核心内容可以概括为：同门课 2022 官方 guest lecture row 明确写出 Thomas Shaf at 3M | MModal，说明这一块是工业 speech system 主题。。
从课程组织上看，这一部分承担的是把局部概念嵌回整条识别流水线中的作用，避免读者把公式或模块孤立地记忆。

记录系统中的附加说明指出：2023 course page exposes only a Guest Lecture placeholder. This chapter is reconstructed from the official 2023 schedule row, the official same-course 2022 guest lecture metadata naming Thomas Shaf at 3M | MModal, and official 3M / Solventum product documentation.

在本章结构中，这一单元被并入“23.1 Guest Lecture Reconstruction: Industrial Speech Systems and Clinical Documentation”，目的是把同一阶段的定义、机制和权衡集中讲清。

\subsection{本章小结}
23.1 Guest Lecture Reconstruction: Industrial Speech Systems and Clinical Documentation 这一部分把 Guest Lecture II 中相邻的教学单元串成了一条连续叙事：既交代了这一阶段要解决的问题，也交代了方法、推理或工程权衡为什么要这样组织。

\section{23.2 Clinical Speech Recognition and Documentation}
根据课程页主题、公开视频字幕以及本讲的 coverage ledger，23.2 Clinical Speech Recognition and Documentation 这一部分主要围绕 Fluency Direct fact sheet 直接描述 front-end speech recognition、clinical documentation、NLU 与 CAPD。 展开。
写作上保持 coverage-first 原则：先把老师在这一阶段真正讲过的问题设定、模型部件、训练或推理逻辑讲清，再压缩成适合复习的结构。

\subsection{Fluency Direct fact sheet 直接描述 front-end speech recognition、clinical documentation、NLU 与 CAPD}
从证据位置看，这一部分对应的课程证据区间是 product fact sheet。课程在这里强调的核心内容可以概括为：Fluency Direct fact sheet 直接描述 front-end speech recognition、clinical documentation、NLU 与 CAPD。。
从课程组织上看，这一部分承担的是把局部概念嵌回整条识别流水线中的作用，避免读者把公式或模块孤立地记忆。

记录系统中的附加说明指出：2023 course page exposes only a Guest Lecture placeholder. This chapter is reconstructed from the official 2023 schedule row, the official same-course 2022 guest lecture metadata naming Thomas Shaf at 3M | MModal, and official 3M / Solventum product documentation.

在本章结构中，这一单元被并入“23.2 Clinical Speech Recognition and Documentation”，目的是把同一阶段的定义、机制和权衡集中讲清。

\subsection{本章小结}
23.2 Clinical Speech Recognition and Documentation 这一部分把 Guest Lecture II 中相邻的教学单元串成了一条连续叙事：既交代了这一阶段要解决的问题，也交代了方法、推理或工程权衡为什么要这样组织。

\section{23.3 Ambient AI and Workflow Integration}
根据课程页主题、公开视频字幕以及本讲的 coverage ledger，23.3 Ambient AI and Workflow Integration 这一部分主要围绕 Fluency Align fact sheet 讨论 ambient clinical documentation、virtual assistant、workflow integration。 展开。
写作上保持 coverage-first 原则：先把老师在这一阶段真正讲过的问题设定、模型部件、训练或推理逻辑讲清，再压缩成适合复习的结构。

\subsection{Fluency Align fact sheet 讨论 ambient clinical documentation、virtual assistant、workflow integration}
从证据位置看，这一部分对应的课程证据区间是 product fact sheet。课程在这里强调的核心内容可以概括为：Fluency Align fact sheet 讨论 ambient clinical documentation、virtual assistant、workflow integration。。
从课程组织上看，这一部分承担的是把局部概念嵌回整条识别流水线中的作用，避免读者把公式或模块孤立地记忆。

记录系统中的附加说明指出：2023 course page exposes only a Guest Lecture placeholder. This chapter is reconstructed from the official 2023 schedule row, the official same-course 2022 guest lecture metadata naming Thomas Shaf at 3M | MModal, and official 3M / Solventum product documentation.

在本章结构中，这一单元被并入“23.3 Ambient AI and Workflow Integration”，目的是把同一阶段的定义、机制和权衡集中讲清。

\subsection{本章小结}
23.3 Ambient AI and Workflow Integration 这一部分把 Guest Lecture II 中相邻的教学单元串成了一条连续叙事：既交代了这一阶段要解决的问题，也交代了方法、推理或工程权衡为什么要这样组织。

\section{23.4 Industrial Deep Learning for Speech Understanding}
根据课程页主题、公开视频字幕以及本讲的 coverage ledger，23.4 Industrial Deep Learning for Speech Understanding 这一部分主要围绕 Fluency Direct deep learning fact sheet 把 speech recognition 与 speech understanding / deep learning 结合起来。 展开。
写作上保持 coverage-first 原则：先把老师在这一阶段真正讲过的问题设定、模型部件、训练或推理逻辑讲清，再压缩成适合复习的结构。

\subsection{Fluency Direct deep learning fact sheet 把 speech recognition 与 speech understanding / deep learning 结合起来}
从证据位置看，这一部分对应的课程证据区间是 product fact sheet。课程在这里强调的核心内容可以概括为：Fluency Direct deep learning fact sheet 把 speech recognition 与 speech understanding / deep learning 结合起来。。
从课程组织上看，这一部分承担的是把局部概念嵌回整条识别流水线中的作用，避免读者把公式或模块孤立地记忆。

记录系统中的附加说明指出：2023 course page exposes only a Guest Lecture placeholder. This chapter is reconstructed from the official 2023 schedule row, the official same-course 2022 guest lecture metadata naming Thomas Shaf at 3M | MModal, and official 3M / Solventum product documentation.

在本章结构中，这一单元被并入“23.4 Industrial Deep Learning for Speech Understanding”，目的是把同一阶段的定义、机制和权衡集中讲清。

\subsection{本章小结}
23.4 Industrial Deep Learning for Speech Understanding 这一部分把 Guest Lecture II 中相邻的教学单元串成了一条连续叙事：既交代了这一阶段要解决的问题，也交代了方法、推理或工程权衡为什么要这样组织。

\section{总结与延伸}
Guest Lecture II 这一讲在整门《Speech Recognition and Understanding》课程中的位置，是把 Guest Lecture II 讲成一个可复用的建模模块。
从教材化角度看，本讲最重要的不是记住单一术语，而是理解它如何嵌入完整的 automatic speech recognition pipeline：输入语音如何被表示、模型如何被训练、解码如何得到最终转写，以及这些设计分别带来什么假设和权衡。
\end{document}
